%% ===========================================================================
%%  frontmatter/lembar-pengesahan.tex — Lembar Pengesahan
%%  Halaman pertama yang diberi angka Romawi: i.
%%  Penanda tangan: dosen pembimbing dan Ketua Program Studi.
%% ===========================================================================
\babfrontmatter{Lembar Pengesahan}

\begin{singlespace}
Laporan Magang dengan judul \emph{\JudulLaporanPDF} ini disusun oleh:

\vspace{1\baselineskip}
\begin{tabular}{@{}l@{\hspace{0.5em}}c@{\hspace{0.5em}}l@{}}
Nama          & : & \NamaPenulis \\
NIM           & : & \NIM \\
Program Studi & : & \ProgramStudi \\
\end{tabular}

\vspace{1\baselineskip}
% TODO: sesuaikan kalimat pengantar pengesahan dengan format prodi.
Telah diperiksa dan disetujui untuk dipertanggungjawabkan sebagai laporan
pelaksanaan magang pada \InstansiMagang{} periode \PeriodeMagang.

\vspace{2\baselineskip}
\begin{center}\Kota, \TanggalPengesahan\end{center}

\vspace{2\baselineskip}
\begin{center}
\begin{minipage}[t]{0.6\textwidth}\centering
  Pembimbing,\\[4\baselineskip]
  \uline{\NamaPembimbing}\\
  \GelarPembimbing\\
  \NIDNPembimbing
\end{minipage}
\end{center}

\vspace{2\baselineskip}
\begin{center}
\begin{minipage}[t]{0.6\textwidth}\centering
  Mengetahui,\\[1\baselineskip]
  Ketua Program Studi,\\[3\baselineskip]
  \uline{\NamaKetuaProgramStudi}\\
  \GelarKetuaProgramStudi\\
  \NIDNKetuaProgramStudi
\end{minipage}
\end{center}
\end{singlespace}
